\documentclass[12pt,a4paper]{article}
\usepackage{ctex}
\usepackage{amsmath,amscd,amsbsy,amssymb,latexsym,url,bm,amsthm}
\usepackage{epsfig,graphicx,subfigure}
\usepackage{enumitem,balance}
\usepackage{wrapfig}
\usepackage{mathrsfs,euscript}
\usepackage[usenames]{xcolor}
\usepackage{hyperref}
\usepackage[vlined,ruled,commentsnumbered,linesnumbered]{algorithm2e}

\newtheorem{theorem}{Theorem}
\newtheorem{lemma}[theorem]{Lemma}
\newtheorem{proposition}[theorem]{Proposition}
\newtheorem{corollary}[theorem]{Corollary}
\newtheorem{exercise}{Exercise}
\newtheorem*{solution}{Solution}
\newtheorem{definition}{Definition}
\theoremstyle{definition}


%\numberwithin{equation}{section}
%\numberwithin{figure}{section}

\renewcommand{\thefootnote}{\fnsymbol{footnote}}

\newcommand{\postscript}[2]
 {\setlength{\epsfxsize}{#2\hsize}
  \centerline{\epsfbox{#1}}}

\renewcommand{\baselinestretch}{1.0}

\setlength{\oddsidemargin}{-0.365in}
\setlength{\evensidemargin}{-0.365in}
\setlength{\topmargin}{-0.3in}
\setlength{\headheight}{0in}
\setlength{\headsep}{0in}
\setlength{\textheight}{10.1in}
\setlength{\textwidth}{7in}
\makeatletter \renewenvironment{proof}[1][Proof] {\par\pushQED{\qed}\normalfont\topsep6\p@\@plus6\p@\relax\trivlist\item[\hskip\labelsep\bfseries#1\@addpunct{.}]\ignorespaces}{\popQED\endtrivlist\@endpefalse} \makeatother
\makeatletter
\renewenvironment{solution}[1][Solution] {\par\pushQED{\qed}\normalfont\topsep6\p@\@plus6\p@\relax\trivlist\item[\hskip\labelsep\bfseries#1\@addpunct{.}]\ignorespaces}{\popQED\endtrivlist\@endpefalse} \makeatother



\begin{document}
\noindent

%========================================================================
\noindent\framebox[\linewidth]{\shortstack[c]{
\Large{\textbf{Exercise Sheet 11}}\vspace{1mm}\\
Discrete Mathematics by Hongfei Fu, 2023.11.20}}

\begin{center}
	\footnotesize{\color{black} \quad Name:\_\_\_\_\_\_\_\_\_  \quad Student ID:\_\_\_\_\_\_\_\_\_ \quad Class: \_\_\_\_\_\_\_\_\_\_\_\_} \\


\end{center}

\begin{enumerate}
\item ([R], Page 650, Exercise 3,5,7,9) For Exercises 3–9, determine whether the graph shown has directed or undirected edges, whether it has multiple edges, and whether it has one or more loops.\\
3. \includegraphics*[scale=0.7]{3@p650.png} Answer:\\
The graph has:
\begin{itemize}
    \item Undirected edges.
    \item No multiple edges.
    \item No loops.
\end{itemize}

\vspace{0.5cm}

5. \includegraphics*[scale=0.7]{5@p650.png} Answer:\\
The graph has:
\begin{itemize}
    \item Undirected edges.
    \item Multiple edges.
    \item Three edges.
\end{itemize}

\vspace{0.5cm}

7. \includegraphics*[scale=0.7]{7@p650.png} Answer:\\
The graph has:
\begin{itemize}
    \item Directed edges.
    \item Multiple edges.
    \item Two loops.
\end{itemize}

\vspace{0.5cm}

9. \includegraphics*[scale=0.7]{9@p650.png} Answer:\\
The graph has:
\begin{itemize}
    \item Directed edges.
    \item Multiple edges.
    \item Two loops.
\end{itemize}

\newpage

\item ([R], Page 665, Exercise 1, 2) In Exercises 1-2 find the degree of each vertex in the given undirected graph. Identify all isolated and pendant vertices.\\
1. \includegraphics*[scale=0.7]{1@p665.png} Answer:\\
\begin{itemize}
    \item deg(a) = 2, deg(b) = 4, deg(c) = 1, deg(d) = 0, deg(e) = 2, deg(f) = 3
    \item isolated vertex: d
    \item pendant vertex: c
\end{itemize}
2. \includegraphics*[scale=0.7]{2@p665.png} Answer:\\
\begin{itemize}
    \item deg(a) = 6, deg(b) = 6, deg(c) = 6, deg(d) = 5, deg(e) = 3
    \item There is no isolated vertex or pendant vertex.
\end{itemize}

\vspace{1cm}

\item ([R], Page 665, Exercise 8) In Exercises 8 find the in-degree and out-degree of each vertex for the given directed multigraph.\\
\includegraphics*[scale=0.7]{8@p665.png} Answer:\\
$deg^{+}(a) = 2, deg^{+}(b) = 4, deg^{+}(c) = 1, deg^{+}(d) = 1$\\
$deg^{-}(a) = 2, deg^{-}(b) = 3, deg^{-}(c) = 2, deg^{-}(d) = 1$
\vspace{1cm}



\end{enumerate}
%========================================================================
\end{document}
